% Inputted by paper_journal/main.tex after the Results section.
% Numbers are traceable to experiments/results/annotation{,_v2,_v3}/ and to
% RESEARCH_LOG.md entries dated 2026-08-16 and 2026-08-17.

\section{Validating the Hallucination Heuristics}
\label{sec:halluc-validation}

Both hallucination measures used in Section~\ref{sec:results} are automatic
lexical heuristics. We report a human validation of each. One survives; the
other does not, and the manner of its failure is itself a result.

A single annotator labelled stratified samples drawn from the held-out
evaluation set, blind to system identity and to the detector's output, which
was revealed only after a row had been answered. Decision rules were fixed in
advance: for the unsupported-claim measure a continuous score was binarised at
$\geq 0.5$, registered before any annotation and never retuned; for the
contradiction measure any non-zero score counted as positive. We emphasise that
this is \emph{human--detector} agreement, not inter-annotator agreement; the
two-annotator study described in our protocol was not performed.

\subsection{The unsupported-claim measure is validated}

On 75 rows spanning all seven systems, the unsupported-claim heuristic reaches
Cohen's $\kappa = 0.749$ (95\% CI $[0.582, 0.890]$) against human judgement,
with 88.0\% observed agreement and no directional skew (McNemar $p = 1.0$;
5 false positives against 4 false negatives). This clears the $\kappa \geq 0.6$
threshold we pre-registered. A sensitivity sweep over binarisation thresholds
from 0.1 to 0.9 peaks at 0.4 ($\kappa = 0.778$) against 0.749 at our registered
cut of 0.5; we report the registered value, since selecting the
$\kappa$-maximising threshold would fit the metric to the reference standard
used to judge it.

The grounding result in Section~\ref{sec:results} therefore rests on a measure
with human validation behind it.

\subsection{The EHR-contradiction measure is invalid}

The same protocol applied to the contradiction heuristic gives
$\kappa = -0.037$ (95\% CI $[-0.101, 0.000]$) on the 49 rows where the measure
is defined --- agreement \emph{below} chance --- with \textbf{zero true
positives} (0 TP, 41 TN, 7 FP, 1 FN). The apparently high raw agreement of
83.7\% is an artefact of a shared negative majority and carries no information
at this prevalence.

The failure mechanism is unambiguous. The detector flags an answer when it
contains a negation cue immediately followed by a term appearing on a diagnosis
or laboratory line of the record. It never checks whether that negation is the
\emph{record's own wording}. ICD-10 descriptions embed negation routinely, and
every one of the seven false positives was triggered this way: \emph{without
myelopathy}, \emph{without complications}, \emph{without hematuria},
\emph{without lesion}, \emph{not carried}. An answer that correctly reproduces
the patient's coded diagnosis is scored as contradicting it. Across the 2{,}025
rows not annotated, 172 of the detector's 179 positives (96\%) fall on the two
question types whose gold answer is a verbatim diagnosis string, confirming the
mechanism is systematic rather than incidental.

We therefore withdraw the EHR-contradiction column and the associated claim of
parity between systems. It compared quantities that do not measure
contradiction.

\subsection{A repaired detector, and the limits of repair}

The premise of the heuristic --- the record asserts $X$, the answer denies $X$
--- is sound; only the second half was ever tested. We repair it by firing only
when the record \emph{positively asserts} the term, so that a term negated
everywhere it appears in the record cannot be contradicted. We additionally
restrict candidate terms to the value side of structured fields and match on
word boundaries. Because the repaired measure detects only negated assertions,
we rename it the \textbf{negation-contradiction rate}.

Its scope is strictly narrower than the original name implied. It captures
\emph{Type~1} errors, where the answer negates something the record asserts. It
cannot capture \emph{Type~2} errors, where the answer asserts something the
record refutes: the single human-flagged contradiction the original detector
missed was of this kind, containing no negation at all. Detecting Type~2
requires clinical reasoning over the record and is out of scope for any lexical
method.

We validated the repair on samples \emph{disjoint} from those that exposed the
defect, since a fix designed in view of particular failures cannot be tested on
them. Table~\ref{tab:detector-validation} reports both rounds. Each carried
constructed contradictions as attention checks, excluded from all statistics;
these were included because a working repair implies a near-uniform negative
response set, which is otherwise indistinguishable from inattentive labelling.
All attention checks were caught in both rounds.

\begin{table}[t]
\centering
\caption{Original versus repaired detector, evaluated on samples disjoint from
those that exposed the defect. Attention checks excluded.}
\label{tab:detector-validation}
\footnotesize
\setlength{\tabcolsep}{4pt}
\begin{tabular}{llccccc}
\toprule
Question type & Detector & TP & FP & FN & Prec. & Rec. \\
\midrule
diagnosis      & original  & 0  & 20 & 0 & 0.00 & --- \\
$(n=30)$       & repaired  & 0  & \textbf{0}  & 0 & --- & --- \\
\midrule
monitoring     & original  & 3  & 0  & 9 & 1.00 & 0.25 \\
$(n=20)$       & repaired  & 11 & 4  & 1 & 0.73 & \textbf{0.92} \\
\bottomrule
\end{tabular}
\end{table}

On diagnosis questions the repair is unambiguous: all 20 false positives are
eliminated and none introduced, a paired McNemar test giving $p = 2.0\times
10^{-6}$. On monitoring questions it \emph{fails} its pre-registered acceptance
bar. We required precision $\geq 0.80$; the observed precision is 0.733 (95\%
CI $[0.449, 0.922]$), short by a single row. We report the failure rather than
relax the bar.

The residual false positives are a single artefact, not four independent ones:
all four arise from the same short template answer, \emph{``No abnormal labs
available''}, in which the negated token \emph{abnormal} is a generic clinical
qualifier rather than a finding, yet appears on a laboratory line of the record.
The correction is evident, but deriving it from these four observations and
validating it on them would repeat the circularity this protocol exists to
avoid, so we leave it as a stated limitation.

\subsection{The headline result: contradiction concentrates where the metric
could not see it}

The comparison on monitoring questions inverts the expected story. There the
original detector was not over-flagging at all --- it was \emph{missing} nine of
twelve genuine contradictions, at a recall of 0.25. The repaired detector
recovers eleven of twelve. Of the twelve rows it newly flags, eight are
confirmed genuine.

The original heuristic therefore fails in \emph{both} directions: it
over-flags on diagnosis questions, where the answer copies coded text
containing negation, and under-detects on monitoring questions, where genuine
contradictions are phrased in ways adjacency matching cannot reach. The
consequence for our results is directional and should be stated plainly: the
contradiction rates reported in Section~\ref{sec:results} are \emph{too low},
not merely noisy.

The concentration is striking. Human annotation confirms twelve genuine
contradictions among twenty monitoring-question rows, against \textbf{zero
among thirty} diagnosis-question rows. Contradiction in this system is not
distributed across question types; it is localised in precisely the type whose
answers are guideline panels rather than record fields --- and it is the type
the original metric was structurally least able to detect. We regard this as
the substantive finding of the validation, and a stronger one than any repaired
rate: an evaluation that reported near-zero contradiction across all systems was
measuring the absence of a lexical pattern, not the absence of contradiction.

\subsection{A measurement caveat for mode-conditional metrics}

One further hazard applies to any per-system rate computed over this measure.
Contradiction is undefined for text-only retrieval, which receives no record to
contradict, so those rows are excluded rather than scored as zero. For systems
that \emph{choose} when to enter that mode, the exclusion is not neutral: on
our data roughly 90\% of all contradictions fall in the quarter of questions
the router sends to text-only retrieval (16 of 18 under $T{+}E$, Fisher
$p = 4.9\times10^{-9}$; 13 of 14 under $T{+}E{+}K$, $p = 6.0\times10^{-8}$).
Computed naively this manufactures a tenfold apparent advantage for the router.
On a matched question set the router and always-on hybrid are indistinguishable.
Any rate reported over a mode-conditional metric must therefore be computed on
a common set of questions, a requirement we did not appreciate in our original
analysis.

\subsection{Summary}

Of our two hallucination measures, one is validated
($\kappa = 0.749$) and one is withdrawn ($\kappa = -0.037$, zero true
positives). The repaired replacement is validated for diagnosis questions and
only partially for monitoring questions, and measures a strictly narrower
phenomenon than its predecessor claimed to. We report this in full because the
negative result is the more useful one: lexical hallucination heuristics of this
family can be wrong in both directions at once, and neither failure is visible
without human adjudication.
